\documentclass{article}

\usepackage{amsmath, amsthm, amssymb, amsfonts}
\usepackage{thmtools}
\usepackage{graphicx}
\usepackage{setspace}
\usepackage{geometry}
\usepackage{float}
\usepackage[backend=biber,style=ieee,sorting=ynt]{biblatex}
\usepackage[hidelinks]{hyperref}
\usepackage[utf8]{inputenc}
\usepackage[spanish]{babel}
\usepackage{framed}
\usepackage[dvipsnames]{xcolor}
\usepackage{tcolorbox}
\usepackage{tikz}
% \usetikzlibrary{positioning}
% \usetikzlibrary{arrows.meta}
\usepackage{caption}
\usepackage{longtable}
\usepackage{pdflscape}
\usepackage{svg}
\usepackage{subcaption}
\usepackage{caption}
\usepackage{multirow}
\usepackage{array}
\usepackage{listings}
\usepackage{cancel}
\usepackage{fancyhdr}
\usepackage{pgfplots}
\pgfplotsset{compat=1.18}


\addbibresource{referencias.bib}


\colorlet{LightGray}{White!90!Periwinkle}
\colorlet{LightOrange}{Orange!15}
\colorlet{LightGreen}{Green!15}



\newcommand{\HRule}[1]{\rule{\linewidth}{#1}}

\declaretheoremstyle[name=Theorem,]{thmsty}
\declaretheorem[style=thmsty,numberwithin=section]{theorem}
\tcolorboxenvironment{theorem}{colback=LightGray}

\declaretheoremstyle[name=Proposition,]{prosty}
\declaretheorem[style=prosty,numberlike=theorem]{proposition}
\tcolorboxenvironment{proposition}{colback=LightOrange}

\declaretheoremstyle[name=Principle,]{prcpsty}
\declaretheorem[style=prcpsty,numberlike=theorem]{principle}
\tcolorboxenvironment{principle}{colback=LightGreen}

\newcolumntype{L}[1]{>{\raggedleft\let\newline\\\arraybackslash\hspace{0pt}}m{#1}}
\newcolumntype{C}[1]{>{\centering\let\newline\\\arraybackslash\hspace{0pt}}m{#1}}
\newcolumntype{R}[1]{>{\raggedright\let\newline\\\arraybackslash\hspace{0pt}}m{#1}}

\setstretch{1.2}
\geometry{
  textheight=9in,
  textwidth=5.5in,
  top=1in,
  headheight=12pt,
  headsep=25pt,
  footskip=30pt
}

\lstdefinestyle{bashstyle}{
  language=bash,
  basicstyle=\ttfamily,
  backgroundcolor=\color{gray!10},
  keywordstyle=\color{blue},
  commentstyle=\color{green!40!black},
  stringstyle=\color{red},
  showstringspaces=false,
  numbers=left,
  numberstyle=\tiny\color{gray},
  breaklines=true,
  breakatwhitespace=true,
  frame=tb,
  rulecolor=\color{black!70},
  framerule=0.5pt,
  tabsize=4,
  captionpos=b
}

\lstdefinestyle{javastyle}{
  language=Java,
  basicstyle=\ttfamily,
  backgroundcolor=\color{gray!10},
  keywordstyle=\color{blue},
  commentstyle=\color{green!40!black},
  stringstyle=\color{red},
  showstringspaces=false,
  numbers=left,
  numberstyle=\tiny\color{gray},
  breaklines=true,
  breakatwhitespace=true,
  frame=tb,
  rulecolor=\color{black!70},
  framerule=0.5pt,
  tabsize=4,
  captionpos=b
}

\lstdefinestyle{pythonstyle}{
  language=Python,
  basicstyle=\ttfamily,
  backgroundcolor=\color{gray!10},
  keywordstyle=\color{blue},
  commentstyle=\color{green!40!black},
  stringstyle=\color{red},
  showstringspaces=false,
  numbers=left,
  numberstyle=\tiny\color{gray},
  breaklines=true,
  breakatwhitespace=true,
  frame=tb,
  rulecolor=\color{black!70},
  framerule=0.5pt,
  tabsize=4,
  captionpos=b,
  morekeywords={ceil,cbrt,fact,sqrt},
}

\lstdefinestyle{cppstyle}{
  language=C++,
  basicstyle=\ttfamily,
  backgroundcolor=\color{gray!10},
  keywordstyle=\color{blue},
  commentstyle=\color{green!40!black},
  stringstyle=\color{red},
  showstringspaces=false,
  numbers=left,
  numberstyle=\tiny\color{gray},
  breaklines=true,
  breakatwhitespace=true,
  frame=tb,
  rulecolor=\color{black!70},
  framerule=0.5pt,
  tabsize=4,
  captionpos=b
}


\setlength{\parindent}{0pt}
\pagestyle{fancy}
\fancyhf{}
\renewcommand{\headrulewidth}{0pt}
\fancyfoot[C]{\thepage}
\fancyfoot[C]{\footnotesize \thepage \quad \textbar{} \quad Este trabajo está bajo una licencia CC BY-SA 4.0.\\ Más info: \url{https://creativecommons.org/licenses/by/4.0/}}


% ------------------------------------------------------------------------------

\begin{document}

% ------------------------------------------------------------------------------
% Cover Page and ToC
% ------------------------------------------------------------------------------

\title{ \normalsize \textsc{}
  \\ [2.0cm]
  \HRule{1.5pt} \\
  \LARGE \textbf{\uppercase{I Hate Mathematics - solución guía}
    \HRule{2.0pt} \\ [0.6cm] \LARGE{Universidad de Bogotá Jorge Tadeo Lozano} \vspace*{13\baselineskip}}
}
\date{}
\author{\textbf{Alvarado Becerra Ludwig - Franco Calderón Alejandro} \\
  Estructuras de datos - 2026-1S}

\maketitle
\thispagestyle{empty}
\newpage

\tableofcontents
\thispagestyle{empty}
\newpage
\setcounter{page}{1}

\section{Interpretación del problema, entrada y salida de datos}


\subsection{Objetivo}

Evaluar múltiples consultas (queries) de fórmulas matemáticas e imprimir sus resultados numéricos. Cada consulta corresponde a una operación específica (suma, multiplicación, potencia, factorial, etc.) con sus respectivos parámetros.

\subsection{Entrada de datos esperada}

El enunciado nos dice de leer un número \(q\) que representa la cantidad de consultas a procesar. Para el ejemplo del enunciado la primera línea es:

\texttt{8}

Este \texttt{8} es \(q\), ahora se deben leer \(q\) consultas. Cada consulta consiste en un string con el nombre de la operación seguido de sus parámetros. Por ejemplo, la segunda línea es:

\texttt{ADD 1 1}

De esta cadena podemos obtener; la operación \texttt{ADD} y sus dos parámetros \texttt{1} y \texttt{1}. Estos datos se pueden almacenar usando la función \lstinline[style=pythonstyle]|split()| que convierte la cadena en una lista.

Para el caso de ejemplo completo tenemos:

\begin{verbatim}
8
ADD 1 1
MUL 2 2
POW 3 3
POLY 2 3
DIST 0 0 3 4
AVG 3 4 8
FACT 5
UNIVERSE 2 2 1
\end{verbatim}

Donde cada línea después de la primera representa una operación diferente con sus respectivos parámetros.

\subsection{Salida de datos esperada}

Para cada consulta, se debe imprimir una línea con el resultado de la operación. Los valores decimales deben ser redondeados a dos decimales usando la función \lstinline[style=pythonstyle]|round()|. Para el ejemplo anterior, la salida esperada es:

\begin{verbatim}
2
4
27
25
25
5
120
796762.03
\end{verbatim} 

\section{Solución conceptual}

\subsection{Implementar funciones básicas}
\label{sec:funciones-basicas}

El problema requiere implementar operaciones matemáticas usando funciones simples. Es importante construir las funciones complejas a partir de las más simples, según indica el enunciado.

\subsubsection{Función ADD}

La función de suma \textbf{ADD}\((a,b)\) retorna \(a + b\). Esta es una operación básica que se implementa directamente con el operador \texttt{+}.

\subsubsection{Función MUL}

La función de multiplicación \textbf{MUL}\((a,b)\) retorna \(a \times b\). Se implementa directamente con el operador \texttt{*}.

\subsubsection{Función POW}

La función de potencia \textbf{POW}\((a,b)\) retorna \(a^b\). El enunciado especifica que NO se debe usar el operador \texttt{**} ni \lstinline[style=pythonstyle]|math.pow()|, por lo tanto se debe implementar usando un ciclo. 

La idea es multiplicar \(a\) consigo mismo \(b\) veces. Se puede usar un ciclo \lstinline[style=pythonstyle]|for| que itere \(b\) veces y en cada iteración multiplique el resultado acumulado por \(a\).

\begin{enumerate}
  \item Inicializar una variable \texttt{resultado} en 1.
  \item Iterar \(b\) veces.
  \item En cada iteración, multiplicar \texttt{resultado} por \(a\).
  \item Retornar \texttt{resultado}.
\end{enumerate}

\subsubsection{Función FACT}

La función factorial \textbf{FACT}\((n)\) retorna \(n!\). El enunciado indica que debe implementarse recursivamente. El factorial se define como:

\[
n! = \begin{cases}
1 & \text{si } n = 0 \\
n \times (n-1)! & \text{si } n > 0
\end{cases}
\]

La implementación recursiva debe tener un caso base cuando \(n = 0\) o \(n = 1\) que retorne 1, y en caso contrario retornar \(n \times \text{FACT}(n-1)\).

\subsubsection{Función POLY}

La función polinómica \textbf{POLY}\((a,b)\) retorna \(a^2 + 2ab + b^2\). Esta fórmula corresponde al desarrollo del binomio \((a+b)^2\). Se puede implementar usando las funciones básicas ya definidas:

\begin{enumerate}
  \item Calcular \(a^2\) usando \texttt{POW(a, 2)}.
  \item Calcular \(2ab\) usando \texttt{MUL(2, MUL(a, b))}.
  \item Calcular \(b^2\) usando \texttt{POW(b, 2)}.
  \item Sumar los tres resultados.
\end{enumerate}

\subsubsection{Función DIST}

La función de distancia \textbf{DIST}\((x_1,y_1,x_2,y_2)\) retorna \((x_1 - x_2)^2 + (y_1 - y_2)^2\). Note que esta NO es la distancia euclidiana completa (que requeriría raíz cuadrada), sino solo la suma de los cuadrados de las diferencias.

\begin{enumerate}
  \item Calcular la diferencia \(x_1 - x_2\).
  \item Elevar al cuadrado usando \texttt{POW}.
  \item Calcular la diferencia \(y_1 - y_2\).
  \item Elevar al cuadrado usando \texttt{POW}.
  \item Sumar ambos resultados.
\end{enumerate}

\subsubsection{Función AVG}

La función de promedio \textbf{AVG}\((a,b,c)\) retorna \(\frac{a+b+c}{3}\). Se suma los tres valores y se divide entre 3.

\subsection{Función SECRET OF UNIVERSE}
\label{sec:universe}

La función más compleja del problema es \textbf{SECRET OF UNIVERSE}\((n,a,b)\) que calcula:

\[
\sum_{i=1}^{n}
\left(
\sum_{j=1}^{n}
\left(
\prod_{k=1}^{n}
\left(\pi \times e^{i + j +k} +  (a^2 + 2ab + b^2) + \frac{a + b + n}{3} \right)
\right)
\right)
\]

Esta fórmula tiene tres niveles de anidación: una suma externa, una suma interna, y un producto más interno. Para implementarla:

\begin{enumerate}
  \item Inicializar una variable para la suma externa.
  \item Iterar \(i\) de 1 a \(n\).
  \item Para cada \(i\), inicializar una variable para la suma interna.
  \item Iterar \(j\) de 1 a \(n\).
  \item Para cada \(j\), inicializar una variable para el producto.
  \item Iterar \(k\) de 1 a \(n\).
  \item Para cada \(k\), calcular el término: \(\pi \times e^{i + j +k} +  (a^2 + 2ab + b^2) + \frac{a + b + n}{3}\)
  \item Multiplicar este término al producto acumulado.
  \item Agregar el producto a la suma interna.
  \item Agregar la suma interna a la suma externa.
  \item Retornar la suma externa.
\end{enumerate}

Para calcular \((a^2 + 2ab + b^2)\) se puede reutilizar la función \texttt{POLY(a, b)}. Para calcular \(\frac{a + b + n}{3}\) se puede usar \texttt{AVG(a, b, n)}. Se pueden usar las constantes \lstinline[style=pythonstyle]|math.pi| y \lstinline[style=pythonstyle]|math.e| según indica el enunciado.

\section{Código modelo}
\label{sec:codigo}

Se va a explicar parte a parte una posible solución al problema, recuerde que puede probar con diferentes implementaciones o incluso formulando más soluciones al problema, se invita al estudiante que haga ese ejercicio.

En primer lugar, se importa el módulo de \lstinline[style=pythonstyle]|math| para utilizar las constantes \lstinline[style=pythonstyle]|math.pi| y \lstinline[style=pythonstyle]|math.e|:

\begin{lstlisting}[style=pythonstyle]
import math
\end{lstlisting}

\subsection{Implementación de funciones básicas}

Siguiendo la teoría de la sección~\ref{sec:funciones-basicas}, se implementan las funciones básicas.

\subsubsection{Funciones ADD y MUL}

Estas funciones son directas, simplemente retornan la operación correspondiente:

\begin{lstlisting}[style=pythonstyle]
def ADD(a: int, b: int) -> int:
    return a + b

def MUL(a: int, b: int) -> int:
    return a * b
\end{lstlisting}

\subsubsection{Función POW}

Para la potencia, se debe usar un ciclo en lugar del operador \texttt{**}:

\begin{lstlisting}[style=pythonstyle]
def POW(a: int, b: int) -> int:
    resultado = 1
    for _ in range(b):
        resultado = MUL(resultado, a)
    return resultado
\end{lstlisting}

Note que se reutiliza la función \texttt{MUL} para la multiplicación.

\subsubsection{Función FACT}

El factorial se implementa de forma recursiva:

\begin{lstlisting}[style=pythonstyle]
def FACT(n: int) -> int:
    if (n == 0 or n == 1):
        return 1
    return MUL(n, FACT(n - 1))
\end{lstlisting}

\subsubsection{Función POLY}

La función polinómica reutiliza las funciones ya definidas:

\begin{lstlisting}[style=pythonstyle]
def POLY(a: int, b: int) -> int:
    a_cuadrado = POW(a, 2)
    dos_ab = MUL(2, MUL(a, b))
    b_cuadrado = POW(b, 2)
    return ADD(ADD(a_cuadrado, dos_ab), b_cuadrado)
\end{lstlisting}

\subsubsection{Función DIST}

Para la distancia al cuadrado:

\begin{lstlisting}[style=pythonstyle]
def DIST(x1: int, y1: int, x2: int, y2: int) -> int:
    diff_x = x1 - x2
    diff_y = y1 - y2
    return ADD(POW(diff_x, 2), POW(diff_y, 2))
\end{lstlisting}

\subsubsection{Función AVG}

El promedio se calcula sumando y dividiendo:

\begin{lstlisting}[style=pythonstyle]
def AVG(a: int, b: int, c: int) -> float:
    suma = ADD(ADD(a, b), c)
    return suma / 3
\end{lstlisting}

\subsection{Función SECRET OF UNIVERSE}

La función más compleja requiere tres ciclos anidados según se explicó en la sección~\ref{sec:universe}:

\begin{lstlisting}[style=pythonstyle]
def UNIVERSE(n: int, a: int, b: int) -> float:
    suma_externa = 0
    poly_valor = POLY(a, b)
    avg_valor = AVG(a, b, n)
    
    for i in range(1, n + 1):
        suma_interna = 0
        for j in range(1, n + 1):
            producto = 1
            for k in range(1, n + 1):
                exponente = i + j + k
                termino = math.pi * (math.e ** exponente) + poly_valor + avg_valor
                producto *= termino
            suma_interna += producto
        suma_externa += suma_interna
    
    return suma_externa
\end{lstlisting}

Note que se calculan \texttt{poly\_valor} y \texttt{avg\_valor} fuera de los ciclos para evitar cálculos repetidos innecesarios.

\subsection{Lectura de datos y procesamiento de consultas}

Ahora se procede a leer el número de consultas:

\begin{lstlisting}[style=pythonstyle]
q = int(input())
\end{lstlisting}

Luego se itera sobre cada consulta:

\begin{lstlisting}[style=pythonstyle]
for _ in range(q):
    data = input().split()
    operacion = data[0]
\end{lstlisting}

La función \lstinline[style=pythonstyle]|split()| convierte la cadena en una lista. Por ejemplo, si el usuario ingresa:

\texttt{"ADD 1 1"}

Después de \lstinline[style=pythonstyle]|split()|, \texttt{data} será:

\begin{lstlisting}[style=pythonstyle]
["ADD", "1", "1"]
\end{lstlisting}

Permitiendo extraer la operación y sus parámetros.

\subsection{Evaluación de operaciones}

Se debe evaluar qué operación ejecutar usando condicionales. Para cada operación se extraen los parámetros necesarios de la lista \texttt{data} y se convierten a enteros:

\begin{lstlisting}[style=pythonstyle]
    if (operacion == "ADD"):
        a, b = int(data[1]), int(data[2])
        resultado = ADD(a, b)
        print(resultado)
        
    elif (operacion == "MUL"):
        a, b = int(data[1]), int(data[2])
        resultado = MUL(a, b)
        print(resultado)
        
    elif (operacion == "POW"):
        a, b = int(data[1]), int(data[2])
        resultado = POW(a, b)
        print(resultado)
        
    elif (operacion == "FACT"):
        n = int(data[1])
        resultado = FACT(n)
        print(resultado)
        
    elif (operacion == "POLY"):
        a, b = int(data[1]), int(data[2])
        resultado = POLY(a, b)
        print(resultado)
        
    elif (operacion == "DIST"):
        x1, y1 = int(data[1]), int(data[2])
        x2, y2 = int(data[3]), int(data[4])
        resultado = DIST(x1, y1, x2, y2)
        print(resultado)
        
    elif (operacion == "AVG"):
        a, b, c = int(data[1]), int(data[2]), int(data[3])
        resultado = AVG(a, b, c)
        print(round(resultado, 2))
        
    elif (operacion == "UNIVERSE"):
        n, a, b = int(data[1]), int(data[2]), int(data[3])
        resultado = UNIVERSE(n, a, b)
        print(round(resultado, 2))
\end{lstlisting}

Note que para \texttt{AVG} y \texttt{UNIVERSE} se usa \lstinline[style=pythonstyle]|round(resultado, 2)| para redondear a dos decimales como lo requiere el enunciado.

\subsection{Código final}

Ensamblando todas las piezas del código obtenemos:

\begin{lstlisting}[style=pythonstyle]
import math

def ADD(a: int, b: int) -> int:
    return a + b

def MUL(a: int, b: int) -> int:
    return a * b

def POW(a: int, b: int) -> int:
    resultado = 1
    for _ in range(b):
        resultado = MUL(resultado, a)
    return resultado

def FACT(n: int) -> int:
    if (n == 0 or n == 1):
        return 1
    return MUL(n, FACT(n - 1))

def POLY(a: int, b: int) -> int:
    a_cuadrado = POW(a, 2)
    dos_ab = MUL(2, MUL(a, b))
    b_cuadrado = POW(b, 2)
    return ADD(ADD(a_cuadrado, dos_ab), b_cuadrado)

def DIST(x1: int, y1: int, x2: int, y2: int) -> int:
    diff_x = x1 - x2
    diff_y = y1 - y2
    return ADD(POW(diff_x, 2), POW(diff_y, 2))

def AVG(a: int, b: int, c: int) -> float:
    suma = ADD(ADD(a, b), c)
    return suma / 3

def UNIVERSE(n: int, a: int, b: int) -> float:
    suma_externa = 0
    poly_valor = POLY(a, b)
    avg_valor = AVG(a, b, n)
    
    for i in range(1, n + 1):
        suma_interna = 0
        for j in range(1, n + 1):
            producto = 1
            for k in range(1, n + 1):
                exponente = i + j + k
                termino = math.pi * (math.e ** exponente) + poly_valor + avg_valor
                producto *= termino
            suma_interna += producto
        suma_externa += suma_interna
    
    return suma_externa

q = int(input())

for _ in range(q):
    data = input().split()
    operacion = data[0]
    
    if (operacion == "ADD"):
        a, b = int(data[1]), int(data[2])
        resultado = ADD(a, b)
        print(resultado)
        
    elif (operacion == "MUL"):
        a, b = int(data[1]), int(data[2])
        resultado = MUL(a, b)
        print(resultado)
        
    elif (operacion == "POW"):
        a, b = int(data[1]), int(data[2])
        resultado = POW(a, b)
        print(resultado)
        
    elif (operacion == "FACT"):
        n = int(data[1])
        resultado = FACT(n)
        print(resultado)
        
    elif (operacion == "POLY"):
        a, b = int(data[1]), int(data[2])
        resultado = POLY(a, b)
        print(resultado)
        
    elif (operacion == "DIST"):
        x1, y1 = int(data[1]), int(data[2])
        x2, y2 = int(data[3]), int(data[4])
        resultado = DIST(x1, y1, x2, y2)
        print(resultado)
        
    elif (operacion == "AVG"):
        a, b, c = int(data[1]), int(data[2]), int(data[3])
        resultado = AVG(a, b, c)
        print(round(resultado, 2))
        
    elif (operacion == "UNIVERSE"):
        n, a, b = int(data[1]), int(data[2]), int(data[3])
        resultado = UNIVERSE(n, a, b)
        print(round(resultado, 2))
\end{lstlisting}



\newpage

\thispagestyle{empty}

\vspace*{0.3\textwidth}

\begin{figure}[H]
  \centering
  \includesvg[width=\textwidth]{img/LogoUtadeo.svg}
\end{figure}



\end{document}
